% This file is generated from the corresponding Obsidian note.
% Source: Teoricos/GS - Teo19.md
\chapter{Corchete de Lie}
\label{ch:teo-19}
\begin{bookdefinition}{Subgrupo monoparamétrico de difeomorfismos}
Una familia de difeomorfismos $\{\theta_t\}_{t\in\mathbb R}$ de $M$ es un subgrupo monoparamétrico de difeomorfismos de $M$ si satisface:

\begin{itemize}
\item \textbf{(i)} $\theta_t:M\to M$ es un difeomorfismo para todo $t\in\mathbb R$.
\item \textbf{(ii)} $\theta_0=\operatorname{Id}_M$.
\item \textbf{(iii)} $\theta_s\circ\theta_t=\theta_{s+t}$, para todo $s,t\in\mathbb R$.
\end{itemize}
Equivalentemente, la aplicación

\[
\xi:(\mathbb R,+)\to(\operatorname{Diff}(M),\circ),\qquad t\mapsto\theta_t,
\]
es un homomorfismo de grupos.

\end{bookdefinition}
\begin{bookproposition}{Generador infinitesimal y su suavidad}
Sea $\theta=\{\theta_t\}_{t\in\mathbb R}$ un subgrupo monoparamétrico suave de difeomorfismos de $M$. Para cada $p\in M$, definimos

\[
X_p(f):=\left.\frac{d}{dt}\right|_{t=0}f(\theta_t(p)),\qquad f\in C^\infty(M).
\]
Entonces $X_p\in T_pM$ para todo $p\in M$, y $X$ es un campo vectorial suave en $M$. Este campo se llama el \textbf{generador infinitesimal} del subgrupo monoparamétrico $\theta$.

\end{bookproposition}
\begin{bookproof}{}
\begin{enumerate}
\item Tomemos una carta suave $(U,\varphi=(x_1,\ldots,x_n))$.
\item Queremos ver que las funciones coordenadas de $X$ en esta carta son suaves.
\item Recordemos que, para $q\in U$,
\end{enumerate}
\[
X_q(x_i)=\left.\frac{d}{dt}\right|_{t=0}x_i(\theta_t(q)).
\]
\begin{enumerate}[start=4]
\item Como $\theta:\mathbb R\times M\to M$ es suave, la función
\end{enumerate}
\[
(t,q)\mapsto x_i(\theta_t(q))
\]
es suave donde está definida. 5. Entonces

\[
X^i(q)=X_q(x_i)=\left.\frac{d}{dt}\right|_{t=0}x_i(\theta_t(q))
\]
es suave como función de $q$. 6. Por lo tanto, en coordenadas

\[
X|_U=\sum_{i=1}^n X^i\frac{\partial}{\partial x_i}
\]
con $X^i\in C^\infty(U)$, y se sigue que $X$ es suave.

\bookqed
\end{bookproof}
\section{Corchete de Lie}
\begin{bookremark}{}
Recordemos que dado $X\in\mathfrak X(M)$ y $f\in C^\infty(M)$, tenemos la función suave

\[
Xf:M\to\mathbb R,\qquad p\mapsto (Xf)(p)=X_p(f).
\]
Osea $Xf\in C^{\infty}(M)$

\end{bookremark}
\begin{bookdefinition}{Corchete de Lie}
Sean $X,Y\in\mathfrak X(M)$. Se define el corchete de Lie de $X$ e $Y$, denotado por $[X,Y]$, como el campo que en $p$ actúa sobre $f\in C^\infty(M)$ de la siguiente manera

\[
[X,Y]_p f=X_p(Yf)-Y_p(Xf).
\]
\end{bookdefinition}
\begin{bookproposition}{Buena definicion corchete de Lie}
El corchete de Lie está bien definido, es decir, $[X,Y]_p\in T_pM$ y ademas $[X,Y]$ es un campo suave.

\end{bookproposition}
\begin{bookproof}{}
\begin{enumerate}
\item Necesitamos ver que $[X,Y]_p$ es vector tangente.
\item \textbf{Linealidad} Sea $f,g\in C^\infty(M)$. Entonces
\end{enumerate}
\[
[X,Y]_p(f+g)=X_p(Y(f+g))-Y_p(X(f+g))=[X,Y]_pf+[X,Y]_pg,
\]
por linealidad de $X_p$ e $Y_p$. Para ver la igualdad intermedia, por ejemplo,

\[
(Y(f+g))(q)=Y_q(f+g)=Y_qf+Y_qg=(Yf)(q)+(Yg)(q),
\]
y por lo tanto $Y(f+g)=Yf+Yg$ como funciones. Análogo con $X(f+g)$. 3. \textbf{Leibniz} Primero por Leibniz

\[
Y(fg)=fYg+gYf.
\]
Esto se prueba facilmente de la misma manera que en 2. 4. Entonces

\[
X(Y(fg))=X(fYg+gYf)=X(fYg)+X(gYf).
\]
\begin{enumerate}[start=5]
\item Usando Leibniz de $X$,
\end{enumerate}
\[
X(fYg)=X(f)Yg+fX(Yg),
\]
y

\[
X(gYf)=X(g)Yf+gX(Yf).
\]
\begin{enumerate}[start=6]
\item Luego
\end{enumerate}
\[
X(Y(fg))=X(f)Yg+fX(Yg)+X(g)Yf+gX(Yf).
\]
\begin{enumerate}[start=7]
\item Del mismo modo,
\end{enumerate}
\[
X(fg)=fXg+gXf,
\]
y por lo tanto

\[
Y(X(fg))=Y(fXg+gXf)=Y(f)Xg+fY(Xg)+Y(g)Xf+gY(Xf).
\]
\begin{enumerate}[start=8]
\item Restando,
\end{enumerate}
\[
[X,Y](fg)=X(Y(fg))-Y(X(fg)).
\]
\begin{enumerate}[start=9]
\item Los términos cruzados se cancelan porque son productos de funciones reales:
\end{enumerate}
\[
X(f)Yg-Y(g)Xf\equiv0,\qquad X(g)Yf-Y(f)Xg\equiv0
\]
cuando se reordenan los productos entre funciones correspondientes. 10. Queda

\[
[X,Y](fg)=f[X,Y](g)+g[X,Y](f).
\]
entonces

\[
[X,Y]_{p}(fg)=f(g)[X,Y]_{p}(g)+g(p)[X,Y]_{p}(f)
\]
\begin{enumerate}[start=11]
\item Finalmente vimos $[X,Y]_p$ es una vector tangente en $p$, es decir $[X,Y]_p\in T_pM$.
\item \textbf{Suavidad} Para esto usemos \hyperref[blk:3accf8]{Caracterizaciones de suavidad} 4.
\item Entonces dada $f\in C^{\infty}(M)$ tenemos que ver que $[X,Y]f=X(Yf)-Y(Xf)$ es suave
\item Primero notamos $Yf,Xf\in C^{\infty}(M)$ por que $X,Y$ son campos suaves entonces $X(Yf),Y(Xf)\in C^{\infty}(M)$ misma razon
\item Luego $[X,Y]f\in C^{\infty}(M)$ por ser resta de cosas en $C^{\infty}(M)$ por lo tanto $[X,Y]$ es campo suave
\end{enumerate}
\bookqed
\end{bookproof}
\begin{bookproposition}{Propiedades del corchete de Lie}
Sean $X,Y,Z\in\mathfrak X(M)$ y $f,g\in C^\infty(M)$. Entonces:

\begin{itemize}
\item \textbf{(i)} $[X,Y]=-[Y,X]$.
\item \textbf{(ii)} Si $(U,\varphi=(x_{1},\ldots,x_{n}))$ entonces
\end{itemize}
\[
\left[\frac{\partial}{\partial x_i},\frac{\partial}{\partial x_j}\right]=0.
\]
\begin{itemize}
\item \textbf{(iii)} $[X,gY]=g[X,Y]+X(g)Y$.
\item \textbf{(iv)} $[fX,gY]=fg[X,Y]+fX(g)Y-gY(f)X$.
\item \textbf{(v)} Identidad de Jacobi:
\end{itemize}
\[
[X,[Y,Z]]+[Y,[Z,X]]+[Z,[X,Y]]=0.
\]
\end{bookproposition}
\begin{bookproof}{}
\begin{itemize}
\item \textbf{(i)}
\end{itemize}
\begin{enumerate}
\item Sale directo:
\end{enumerate}
\[
-[Y,X]f=-Y(Xf)+X(Yf)=[X,Y]f.
\]
\begin{itemize}
\item \textbf{(ii)}
\end{itemize}
\begin{enumerate}
\item Tomemos $f\in C^\infty(U)$.
\end{enumerate}
\[
\begin{aligned}\left.\frac{\partial}{\partial x_i}\right|_p\left(\frac{\partial f}{\partial x_j}\right)&=\frac{\partial}{\partial r_i}\left[\left(\frac{\partial f}{\partial x_j}\right)\circ\varphi^{-1}\right](\varphi(p))\\&=\frac{\partial}{\partial r_i}\left[\left(\left(\frac{\partial}{\partial r_j}(f\circ\varphi^{-1})\right)\circ\varphi\right)\circ\varphi^{-1}\right](\varphi(p))\\&=\frac{\partial}{\partial r_i}\left[\frac{\partial}{\partial r_j}(f\circ\varphi^{-1})\circ(\varphi\circ\varphi^{-1})\right](\varphi(p))\\&=\frac{\partial}{\partial r_i}\left[\frac{\partial}{\partial r_j}(f\circ\varphi^{-1})\right](\varphi(p))\\&=\frac{\partial^2(f\circ\varphi^{-1})}{\partial r_i\partial r_j}(\varphi(p)).\end{aligned}
\]
\begin{enumerate}[start=2]
\item Y se termina con el teorema de igualdad de las derivadas cruzadas.
\end{enumerate}
\begin{itemize}
\item \textbf{(iii)}
\end{itemize}
\begin{enumerate}
\item Calculamos:
\end{enumerate}
\[
[X,gY]h=X(gYh)-gY(Xh).
\]
\begin{enumerate}[start=2]
\item Usando Leibniz:
\end{enumerate}
\[
X(gYh)=X(g)Yh+gX(Yh).
\]
\begin{enumerate}[start=3]
\item Entonces
\end{enumerate}
\[
[X,gY]h=X(g)Yh+gX(Yh)-gY(Xh)=X(g)Yh+g[X,Y]h.
\]
\begin{enumerate}[start=4]
\item Por lo tanto,
\end{enumerate}
\[
[X,gY]=g[X,Y]+X(g)Y.
\]
\begin{itemize}
\item \textbf{(iv)}
\end{itemize}
\begin{enumerate}
\item Usando \textbf{(iii)} y antisimetría:
\end{enumerate}
\[
[fX,gY]=f[X,gY]-gY(f)X.
\]
\begin{enumerate}[start=2]
\item Luego
\end{enumerate}
\[
[fX,gY]=fg[X,Y]+fX(g)Y-gY(f)X.
\]
\begin{itemize}
\item \textbf{(v)}
\end{itemize}
\begin{enumerate}
\item La prueba de Jacobi es un cálculo. Para $f\in C^\infty(M)$:
\end{enumerate}
\[
\begin{aligned} \big[X,[Y,Z]]f+[Y,[Z,X]]f+[Z,[X,Y]]f&=X[Y,Z]f-[Y,Z]Xf+Y[Z,X]f-[Z,X]Yf+Z[X,Y]f-[X,Y]Zf\\&=XYZf-XZYf-YZXf+ZYXf+YZXf-YXZf\\&\quad -ZXYf+XZYf+ZXYf-ZYXf-XYZf+YXZf\\&=0.\end{aligned}
\]
\bookqed
\end{bookproof}
\phantomsection\label{blk:411d32}
\begin{bookdefinition}{Campos $F$-relacionados}
Sean $F:M\to N$ suave, $X\in\mathfrak X(M)$ y $\widetilde X\in\mathfrak X(N)$. Se dice que $X$ y $\widetilde X$ están $F$-relacionados si

\[
(dF)_pX_p=\widetilde X_{F(p)},\qquad p\in M.
\]
\end{bookdefinition}
\begin{bookproposition}{Naturalidad del corchete}
Sean $F:M\to N$ suave, $X,Y\in\mathfrak X(M)$ y $\widetilde X,\widetilde Y\in\mathfrak X(N)$ tales que $X,Y$ estan $F$-relacionados con $\widetilde{X},\widetilde Y$ respectivamente entonces $[X,Y]$ está $F$-relacionado con $[\widetilde X,\widetilde Y]$.

\end{bookproposition}
\begin{bookproof}{}
\begin{enumerate}
\item Primero notamos
\end{enumerate}
\[
(dF)_{p}[X,Y]_{p}f=[X,Y]_{p}(f\circ F)=X_{p}(Y(f\circ F))-Y_{p}(X(f\circ F))
\]
\begin{enumerate}[start=2]
\item Tenemos que $(dF)_{p}(X_{p})=\widetilde{X}_{F(p)}$ entonces
\end{enumerate}
\[
\widetilde X_{F(p)}(f)=(\widetilde Xf)(F(p)).
\]
Por lo tanto

\[
X_p(f\circ F)=((\widetilde Xf)\circ F)(p).
\]
Como esto vale para todo $p\in M$, se obtiene la igualdad de funciones

\[
X(f\circ F)=(\widetilde Xf)\circ F.
\]
Analogo con $Y(f\circ F)$. 3. Por lo tanto

\[
\begin{aligned}(dF)_{p}[X,Y]_{p}h=[X,Y]_{p}(h\circ F)&=X_{p}((\widetilde Y f)\circ F)-Y_{p}((\widetilde X f)\circ F)\\&=(dF)_{p}(X_{p})(\widetilde{Y}f)-(dF)_{p}(Y_{p})(\widetilde{X}f)\\&=\widetilde{X}_{F(p)}(\widetilde{Y}f)-\widetilde{Y}_{F(p)}(\widetilde{X}f)\\&=[\widetilde{X},\widetilde{Y}]_{F(p)}f\end{aligned}
\]
mostrando que $[X,Y]$ y $[\widetilde{X},\widetilde{Y}]$ estan $F$-relacionados

\bookqed
\end{bookproof}
\section{Derivada de Lie}
\begin{booklemma}{}
Sea $\alpha:(-\epsilon,\epsilon)\to TM$ una curva suave. Entonces, para toda $f\in C^\infty(M)$, se cumple

\[
\alpha'(0)(f)=\frac{d}{dt}\bigg|_{t=0}\alpha(t)f.
\]
donde identificamos $T_{\alpha(0)}(TM)$ con $T_pM$.

\end{booklemma}
\begin{bookproof}{}
\begin{enumerate}
\item Fijemos una carta suave $(U,\varphi=(x_1,\ldots,x_n))$ en $p$.
\item Luego
\end{enumerate}
\[
\alpha(t)=\sum_{i=1}^n b^i(t)\frac{\partial}{\partial x_i}\bigg|_p,
\]
con $b^i:(-\epsilon,\epsilon)\to\mathbb R$ suaves. 3. Vemos a $\alpha'(0)$ como elemento de $T_pM$ vía la identificación $T_{\alpha(0)}(T_pM)\cong T_pM$. 4. Usando la base $\left\{\frac{\partial}{\partial x_1}\bigg|_p,\ldots,\frac{\partial}{\partial x_n}\bigg|_p\right\}$, tenemos

\[
\alpha'(0)=\sum_{i=1}^n (b^i)'(0)\frac{\partial}{\partial x_i}\bigg|_p.
\]
\begin{enumerate}[start=5]
\item Así,
\end{enumerate}
\[
\alpha'(0)f=\sum_{i=1}^n (b^i)'(0)\frac{\partial f}{\partial x_i}(p).
\]
\begin{enumerate}[start=6]
\item Por otro lado,
\end{enumerate}
\[
\alpha(t)f=\sum_{i=1}^n b^i(t)\frac{\partial f}{\partial x_i}(p).
\]
\begin{enumerate}[start=7]
\item Entonces,
\end{enumerate}
\[
\frac{d}{dt}\bigg|_{t=0}\alpha(t)f=\frac{d}{dt}\bigg|_{t=0}\sum_{i=1}^n b^i(t)\frac{\partial f}{\partial x_i}(p)=\sum_{i=1}^n (b^i)'(0)\frac{\partial f}{\partial x_i}(p).
\]
\begin{enumerate}[start=8]
\item Por lo tanto,
\end{enumerate}
\[
\alpha'(0)(f)=\frac{d}{dt}\bigg|_{t=0}\alpha(t)f.
\]
\bookqed
\end{bookproof}
\begin{bookexercise}{Ejercicio 23 del practico}
Sean $M$ una variedad diferenciable, $X,Y\in\mathfrak X(M)$ y sea $p\in M$. Probar que

\[
t\mapsto(d\theta_{-t})_{\theta_t(p)}(Y_{\theta_t(p)})
\]
define una curva $C^\infty$ en $T_pM$, donde $\theta_t$ es el flujo local asociado a $X$.

\end{bookexercise}
\begin{bookproof}{}
\hyperref[blk:d6291e]{Ejercicio}

\bookqed
\end{bookproof}
\begin{booktheorem}{Derivada de Lie NO FINAL}
Para todo $X,Y\in\mathfrak X(M)$ y para todo $p\in M$, se cumple

\[
[X,Y]_p=(\mathcal L_XY)_p=\lim_{t\to0}\frac{(d\theta_{-t})_{\theta_t(p)}Y_{\theta_t(p)}-Y_p}{t},
\]
donde $\theta$ es el flujo de $X$.

\end{booktheorem}
\begin{bookproof}{}
9. Se calculan ambos lados y se prueba que son iguales. 10. Supongamos primero que $X_p\neq0$. 11. Por el teorema de rectificación de campos, existe una carta suave $(U,\varphi=(x_1,\ldots,x_n))$ centrada en $p$ tal que

\[
X=\frac{\partial}{\partial x_1}.
\]
\begin{enumerate}[start=12]
\item Escribimos
\end{enumerate}
\[
Y=\sum_{j=1}^n a_j\frac{\partial}{\partial x_j},
\]
con $a_1,\ldots,a_n$ funciones suaves sobre $U$. 13. Entonces

\[
[X,Y]=\left[\frac{\partial}{\partial x_1},\sum_{j=1}^n a_j\frac{\partial}{\partial x_j}\right]=\sum_{j=1}^n a_j\left[\frac{\partial}{\partial x_1},\frac{\partial}{\partial x_j}\right]+\sum_{j=1}^n \frac{\partial a_j}{\partial x_1}\frac{\partial}{\partial x_j}.
\]
\begin{enumerate}[start=14]
\item Como los campos coordenados conmutan,
\end{enumerate}
\[
\left[\frac{\partial}{\partial x_1},\frac{\partial}{\partial x_j}\right]=0.
\]
\begin{enumerate}[start=15]
\item Por lo tanto,
\end{enumerate}
\[
[X,Y]_p=\sum_{j=1}^n \frac{\partial a_j}{\partial x_1}(p)\frac{\partial}{\partial x_j}\bigg|_p.
\]
\begin{enumerate}[start=16]
\item Por otro lado, consideremos la curva suave
\end{enumerate}
\[
\alpha(t)=(d\theta_{-t})_{\theta_t(p)}Y_{\theta_t(p)}\in T_pM.
\]
\begin{enumerate}[start=17]
\item Queremos calcular
\end{enumerate}
\[
\alpha'(0)\in T_{\alpha(0)}(T_pM)\simeq T_pM.
\]
\begin{enumerate}[start=18]
\item Entonces
\end{enumerate}
\[
\alpha'(0)=\sum_{j=1}^n c_j\frac{\partial}{\partial x_j}\bigg|_p,
\]
donde

\[
c_j=\alpha'(0)(x_j).
\]
\begin{enumerate}[start=19]
\item Por el lema previo,
\end{enumerate}
\[
\alpha'(0)(x_j)=\frac{d}{dt}\bigg|_{t=0}\alpha(t)(x_j).
\]
\begin{enumerate}[start=20]
\item Luego
\end{enumerate}
\[
\alpha'(0)(x_j)=\frac{d}{dt}\bigg|_{t=0}\left((d\theta_{-t})_{\theta_t(p)}Y_{\theta_t(p)}\right)(x_j).
\]
\begin{enumerate}[start=21]
\item Por definición de diferencial,
\end{enumerate}
\[
\left((d\theta_{-t})_{\theta_t(p)}Y_{\theta_t(p)}\right)(x_j)=Y_{\theta_t(p)}(x_j\circ\theta_{-t}).
\]
\begin{enumerate}[start=22]
\item Como
\end{enumerate}
\[
Y_{\theta_t(p)}=\sum_{i=1}^n a_i(\theta_t(p))\frac{\partial}{\partial x_i}\bigg|_{\theta_t(p)},
\]
tenemos

\[
Y_{\theta_t(p)}(x_j\circ\theta_{-t})=\sum_{i=1}^n a_i(\theta_t(p))\frac{\partial}{\partial x_i}\bigg|_{\theta_t(p)}(x_j\circ\theta_{-t}).
\]
\begin{enumerate}[start=23]
\item En estas coordenadas, como $X=\frac{\partial}{\partial x_1}$, el flujo satisface
\end{enumerate}
\[
x_j\circ\theta_{-t}=x_j-t\delta_{1j}.
\]
\begin{enumerate}[start=24]
\item Por lo tanto,
\end{enumerate}
\[
\frac{\partial}{\partial x_i}\bigg|_{\theta_t(p)}(x_j\circ\theta_{-t})=\delta_{ij}.
\]
\begin{enumerate}[start=25]
\item Así,
\end{enumerate}
\[
Y_{\theta_t(p)}(x_j\circ\theta_{-t})=a_j(\theta_t(p)).
\]
\begin{enumerate}[start=26]
\item Entonces
\end{enumerate}
\[
\alpha'(0)(x_j)=\frac{d}{dt}\bigg|_{t=0}a_j(\theta_t(p))=X_pa_j=\frac{\partial a_j}{\partial x_1}(p).
\]
\begin{enumerate}[start=27]
\item El igual del medio se justifica así: si $v\in T_pM$ y $\gamma$ es una curva con $\gamma(0)=p$ y $\gamma'(0)=v$, entonces
\end{enumerate}
\[
v(f)=\frac{d}{dt}\bigg|_{t=0}f(\gamma(t)).
\]
\begin{enumerate}[start=28]
\item Por lo tanto,
\end{enumerate}
\[
\alpha'(0)=\sum_{j=1}^n \frac{\partial a_j}{\partial x_1}(p)\frac{\partial}{\partial x_j}\bigg|_p=[X,Y]_p.
\]
\begin{enumerate}[start=29]
\item Ahora supongamos que $X_p=0$.
\item En este caso, la curva integral de $X$ que parte de $p$ es constante, luego
\end{enumerate}
\[
\theta_t(p)=p.
\]
\begin{enumerate}[start=31]
\item Entonces la curva queda
\end{enumerate}
\[
\alpha(t)=(d\theta_{-t})_pY_p.
\]
\begin{enumerate}[start=32]
\item Para toda $f\in C^\infty(M)$,
\end{enumerate}
\[
\alpha'(0)(f)=\frac{d}{dt}\bigg|_{t=0}\alpha(t)(f)=\frac{d}{dt}\bigg|_{t=0}Y_p(f\circ\theta_{-t}).
\]
\begin{enumerate}[start=33]
\item Como $Y_p$ es una derivación fija,
\end{enumerate}
\[
\alpha'(0)(f)=Y_p\left(\frac{d}{dt}\bigg|_{t=0}f\circ\theta_{-t}\right).
\]
\begin{enumerate}[start=34]
\item Pero
\end{enumerate}
\[
\frac{d}{dt}\bigg|_{t=0}f(\theta_{-t}(q))=-X_qf.
\]
\begin{enumerate}[start=35]
\item Es decir,
\end{enumerate}
\[
\frac{d}{dt}\bigg|_{t=0}f\circ\theta_{-t}=-Xf.
\]
\begin{enumerate}[start=36]
\item Por lo tanto,
\end{enumerate}
\[
\alpha'(0)(f)=-Y_p(Xf).
\]
\begin{enumerate}[start=37]
\item Como $X_p=0$, también
\end{enumerate}
\[
[X,Y]_pf=X_p(Yf)-Y_p(Xf)=-Y_p(Xf).
\]
\begin{enumerate}[start=38]
\item Así,
\end{enumerate}
\[
\alpha'(0)(f)=[X,Y]_pf.
\]
\begin{enumerate}[start=39]
\item Como esto vale para toda $f\in C^\infty(M)$, concluimos que
\end{enumerate}
\[
\alpha'(0)=[X,Y]_p.
\]
\begin{enumerate}[start=40]
\item Por lo tanto,
\end{enumerate}
\[
(\mathcal L_XY)_p=[X,Y]_p.
\]
\bookqed
\end{bookproof}
\begin{bookdefinition}{Corchete de Lie mediante el flujo}
Sean $X,Y\in\mathfrak X(M)$, y sea $\theta$ el flujo local de $X$. Para $p\in M$, consideramos la curva

\[
t\mapsto(d\theta_{-t})_{\theta_t(p)}Y_{\theta_t(p)}
\]
en $T_pM$. Su derivada en $t=0$ es el corchete:

\[
[X,Y]_p=\left.\frac{d}{dt}\right|_{0}(d\theta_{-t})_{\theta_t(p)}Y_{\theta_t(p)}.
\]
Esto se puede pensar como que empujamos $Y$ por el flujo de $X$ y vemos cómo cambia infinitesimalmente.

\end{bookdefinition}
\section{Lema de uniformidad conjunta para dos flujos}
\begin{booklemma}{Uniformidad conjunta para dos flujos}
Sean $X$ e $Y$ dos campos en una variedad $M$, y sean $\theta$ y $\phi$ sus respectivos flujos. Para todo $p\in M$ existen un entorno abierto $U$ en $M$ de $p$ y $\delta>0$ tal que $\theta_s\circ\phi_t$ y $\phi_t\circ\theta_s$ están definidos en $U$ para $|s|<\delta$ y $|t|<\delta$.

\end{booklemma}
\begin{bookproof}{}
\begin{enumerate}
\item Usemos uniformidad para cada uno de $X$ e $Y$.
\item Existen abiertos $V_0$ y $W_0$ de $M$ en $p$, y $\varepsilon_1,\varepsilon_2>0$, tales que
\end{enumerate}
\[
(-\varepsilon_1,\varepsilon_1)\times V_0\subset\mathcal D_\theta,
\]
\[
(-\varepsilon_2,\varepsilon_2)\times W_0\subset\mathcal D_\phi.
\]
\begin{enumerate}[start=3]
\item Por tanto están bien definidos
\end{enumerate}
\[
\theta:(-\varepsilon_1,\varepsilon_1)\times V_0\to M,
\]
\[
\phi:(-\varepsilon_2,\varepsilon_2)\times W_0\to M.
\]
\begin{enumerate}[start=4]
\item Después vamos a querer hacer $\theta_s\circ\phi_t$ en cierto abierto de $p$. En particular $\phi_t(q)$ debe poder evaluarse en $\theta_s$.
\item Como también $\theta_s(q)$ debe poder evaluarse en $\phi_t$, pensamos en los abiertos de $\mathbb R\times M$:
\end{enumerate}
\[
\theta^{-1}(V_0),\qquad \phi^{-1}(W_0).
\]
\begin{enumerate}[start=6]
\item Como $\mathbb R\times M$ tiene la topología producto, existe $\delta>0$ y un abierto $U$ de $p$ tal que
\end{enumerate}
\[
(-\delta,\delta)\times U\subset \theta^{-1}(W_0)\cap\phi^{-1}(V_0),
\]
con $\delta<\varepsilon_1,\varepsilon_2$. 7. Resta ver que $\delta$ y $U$ cumplen lo que queremos. 8. Sean $s,t\in(-\delta,\delta)$ y $q\in U$. 9. Tenemos $(t,q)\in(-\varepsilon_2,\varepsilon_2)\times W_0$, por tanto $\phi_t(q)$ está definido. 10. Además, como $(t,q)\in\phi^{-1}(V_0)$, se tiene $\phi_t(q)\in V_0$. 11. Como $s\in(-\varepsilon_1,\varepsilon_1)$ y $\phi_t(q)\in V_0$, podemos evaluar $\theta_s(\phi_t(q))$. 12. Similarmente, $(s,q)\in\theta^{-1}(W_0)$ implica $\theta_s(q)\in W_0$, y entonces $\phi_t(\theta_s(q))$ también está definido.

\bookqed
\end{bookproof}
\phantomsection\label{blk:39cc9a}
\begin{bookproposition}{Criterio de conmutación de flujos (Hulet 9.15)}
Sean $X,Y\in\mathfrak X(M)$ tales que $[X,Y]=0$ y sean $\theta,\phi$ los flujos respectivos. Dado $p\in M$, sean $U$ y $\delta$ los del \hyperref[blk:39cc9a]{Uniformidad conjunta para dos flujos}. Entonces:

\begin{itemize}
\item \textbf{(i)} $(d\theta_s)_qY_q=Y_{\theta_s(q)}$ con $q\in U$ para $s\in(-\delta,\delta)$.
\item \textbf{(ii)} $\theta_t\circ\phi_s=\phi_s\circ\theta_t$ en $U$ para $t,s\in(-\delta,\delta)$.
\end{itemize}
\end{bookproposition}
\begin{bookproof}{}
\begin{itemize}
\item (i)
\end{itemize}
\begin{enumerate}
\item Recordemos la fórmula
\end{enumerate}
\[
[X,Y](p)=\frac{d}{dt}\Big|_{t=0}(d\theta_{-t})_{\theta_t(p)}\big(Y_{\theta_t(p)}\big),\qquad p\in M.
\]
\begin{enumerate}[start=2]
\item Notar que probar $(i)$ equivale a probar que
\end{enumerate}
\[
Y_p=(d\theta_{-s})_{\theta_s(p)}\big(Y_{\theta_s(p)}\big),\qquad \forall s\in (-\delta ,\delta),\forall p\in U
\]
\begin{enumerate}[start=3]
\item En otras palabras, para probar $(i)$ debemos ver que, para cada $p$, la curva
\end{enumerate}
\[
z(s)=(d\theta_{-s})_{\theta_s(p)}\big(Y_{\theta_s(p)}\big)
\]
es constantemente igual a $Y_p$. 4. Calculamos $z'(s)$. Por definición,

\[
z'(s)=\frac{d}{dt}\Big|_{t=0}z(t+s).
\]
\begin{enumerate}[start=5]
\item Ahora,
\end{enumerate}
\[
z(t+s)=(d\theta_{-(t+s)})_{\theta_{t+s}(p)}\big(Y_{\theta_{t+s}(p)}\big).
\]
\begin{enumerate}[start=6]
\item Como
\end{enumerate}
\[
\theta_{t+s}(p)=\theta_t(\theta_s(p))
\]
y

\[
\theta_{-(t+s)}=\theta_{-s}\circ\theta_{-t},
\]
por la regla de la cadena obtenemos

\[
z(t+s)=(d\theta_{-s})_{\theta_s(p)}\left((d\theta_{-t})_{\theta_t(\theta_s(p))}\big(Y_{\theta_t(\theta_s(p))}\big)\right).
\]
\begin{enumerate}[start=7]
\item Por lo tanto,
\end{enumerate}
\[
z'(s)=\frac{d}{dt}\Big|_{t=0}(d\theta_{-s})_{\theta_s(p)}\left((d\theta_{-t})_{\theta_t(\theta_s(p))}\big(Y_{\theta_t(\theta_s(p))}\big)\right).
\]
\begin{enumerate}[start=8]
\item Como $(d\theta_{-s})_{\theta_s(p)}$ es lineal y no depende de $t$,
\end{enumerate}
\[
z'(s)=(d\theta_{-s})_{\theta_s(p)}\left(\frac{d}{dt}\Big|_{t=0}(d\theta_{-t})_{\theta_t(\theta_s(p))}\big(Y_{\theta_t(\theta_s(p))}\big)\right).
\]
\begin{enumerate}[start=9]
\item Por la fórmula del corchete aplicada al punto $\theta_s(p)$,
\end{enumerate}
\[
\frac{d}{dt}\Big|_{t=0}(d\theta_{-t})_{\theta_t(\theta_s(p))}\big(Y_{\theta_t(\theta_s(p))}\big)=[X,Y]_{\theta_s(p)}.
\]
\begin{enumerate}[start=10]
\item Luego
\end{enumerate}
\[
z'(s)=(d\theta_{-s})_{\theta_s(p)}\big([X,Y]_{\theta_s(p)}\big).
\]
\begin{enumerate}[start=11]
\item Como $[X,Y]=0$, concluimos que
\end{enumerate}
\[
z'(s)=0,\qquad \forall s\in (-\delta ,\delta )
\]
osea $z(s)$ es constante 12. Además,

\[
z(0)=(d\theta_0)_p(Y_p)=Y_p.
\]
\begin{enumerate}[start=13]
\item Por lo tanto,
\end{enumerate}
\[
z(s)=Y_p,\qquad \forall s,
\]
es decir,

\[
(d\theta_{-t})_{\theta_t(p)}\big(Y_{\theta_t(p)}\big)=Y_p.
\]
\begin{enumerate}[start=14]
\item Aplicando $(d\theta_t)_p$ a ambos lados, obtenemos
\end{enumerate}
\[
(d\theta_t)_p(Y_p)=Y_{\theta_t(p)}.
\]
\begin{enumerate}[start=15]
\item \textbf{(ii)} Ahora fijemos $t$ y $q$. Consideramos
\end{enumerate}
\[
\alpha(s)=\theta_t(\phi_s(q)),
\]
\[
\beta(s)=\phi_s(\theta_t(q)).
\]
\begin{enumerate}[start=16]
\item Tenemos
\end{enumerate}
\[
\alpha(0)=\theta_t(q)=\beta(0).
\]
\begin{enumerate}[start=17]
\item Además,
\end{enumerate}
\[
\alpha'(s)=(d\theta_t)_{\phi_s(q)}Y_{\phi_s(q)}=Y_{\theta_t(\phi_s(q))}=Y_{\alpha(s)},
\]
donde usamos \textbf{(i)}. 18. Por otro lado,

\[
\beta'(s)=Y_{\phi_s(\theta_t(q))}=Y_{\beta(s)}.
\]
\begin{enumerate}[start=19]
\item Entonces $\alpha$ y $\beta$ son curvas integrales de $Y$ con el mismo dato inicial.
\item Por unicidad,
\end{enumerate}
\[
\theta_t(\phi_s(q))=\phi_s(\theta_t(q)).
\]
\bookqed
\end{bookproof}
\phantomsection\label{blk:26db5d}
\begin{bookproposition}{(Hulet 9.16)}
Sean $X,Y\in\mathfrak{X}(M)$, cuyos flujos conmutan. Entonces

\[
[X,Y]=0.
\]
\end{bookproposition}
\begin{bookproof}{}
\begin{enumerate}
\item Sean $\theta$ y $\phi$ los flujos de $X$ e $Y$, respectivamente. Fijemos $p\in M$. Por el lema de uniformidad conjunta de los flujos, existen un abierto $U$ que contiene a $p$ y $\delta>0$ tales que, para $|s|,|t|<\delta$, las composiciones $\theta_t\circ\phi_s$ y $\phi_s\circ\theta_t$ están definidas en $U$.
\item Como los flujos conmutan, para $|s|,|t|<\delta$ se cumple
\end{enumerate}
\[
\theta_t(\phi_s(p))=\phi_s(\theta_t(p)).
\]
\begin{enumerate}[start=3]
\item Derivando respecto de $s$ en $s=0$, obtenemos
\end{enumerate}
\[
(d\theta_t)_p(Y_p)=(d\theta_t)_{\phi_0(p)}\left(\frac{d}{ds}\bigg|_{s=0}\phi_s(p)\right)=\frac{d}{ds}\bigg|_{s=0}\theta_t(\phi_s(p))=\frac{d}{ds}\bigg|_{s=0}\phi_s(\theta_t(p))=Y_{\theta_t(p)}.
\]
\begin{enumerate}[start=4]
\item Aplicamos $(d\theta_{-t})_{\theta_t(p)}$ a ambos lados. Entonces
\end{enumerate}
\[
\begin{aligned}(d\theta_{-t})_{\theta_t(p)}\left((d\theta_t)_p(Y_p)\right)&=(d\theta_{-t})_{\theta_t(p)}\left(Y_{\theta_t(p)}\right),\\d(\theta_{-t}\circ\theta_t)_p(Y_p)&=(d\theta_{-t})_{\theta_t(p)}\left(Y_{\theta_t(p)}\right),\\Y_p&=(d\theta_{-t})_{\theta_t(p)}\left(Y_{\theta_t(p)}\right).\end{aligned}
\]
\begin{enumerate}[start=5]
\item Por lo tanto, la curva
\end{enumerate}
\[
t\mapsto(d\theta_{-t})_{\theta_t(p)}\left(Y_{\theta_t(p)}\right)
\]
es constante para $|t|<\delta$. 6. Usando la fórmula del corchete en términos del flujo de $X$,

\[
[X,Y]_p=\frac{d}{dt}\bigg|_{t=0}(d\theta_{-t})_{\theta_t(p)}\left(Y_{\theta_t(p)}\right)=0.
\]
\begin{enumerate}[start=7]
\item Como $p\in M$ era arbitrario, concluimos que
\end{enumerate}
\[
[X,Y]=0.
\]
\bookqed
\end{bookproof}
\begin{bookremark}{}
Recordar que si $\frac{\partial}{\partial x_i}$ y $\frac{\partial}{\partial x_j}$ son los campos de un marco coordenado, entonces

\[
\left[\frac{\partial}{\partial x_i},\frac{\partial}{\partial x_j}\right]=0.
\]
Por tanto, si queremos rectificar dos campos $X,Y$ necesitamos como hipótesis $[X,Y]=0$. Por que si pudieramos rectificarlos entonces

\[
[X,Y]=\left[\frac{\partial}{\partial x_i},\frac{\partial}{\partial x_j}\right]=0
\]
\end{bookremark}
\begin{bookproposition}{Rectificación simultánea (Hulet 9.18)}
Sean $X_1,\ldots,X_k\in\mathfrak X(M)$ linealmente independientes en un abierto de $M$ en $p$ tales que

\[
[X_{i},X_{j}]=0\quad 1\leq i,j\leq k
\]
Entonces existe un sistema coordenado $(U,\varphi=(x_1,\ldots,x_n))$ centrado en $p$ tal que

\[
X_1|_U=\frac{\partial}{\partial x_1},\ldots,\qquad X_k|_U=\frac{\partial}{\partial x_k}
\]
\end{bookproposition}
\begin{bookproof}{}
\begin{enumerate}
\item Por \hyperref[blk:39cc9a]{Uniformidad conjunta para dos flujos}. Existe $\varepsilon>0$ y un abierto $V$ de $p$ tal que la composición
\end{enumerate}
\[
\theta^{1}_{t_1}\circ\cdots\circ\theta^k_{t_k}
\]
está bien definida sobre $V$ para todos $t_1,\ldots,t_k\in(-\varepsilon,\varepsilon)$. 2. Como los campos son linealmente independientes. Achicando $V$ si hace falta, tomamos una carta cúbica centrada en $p$,

\[
(V,\psi=(x_1,\ldots,x_n)),
\]
tal que

\[
\left.\frac{\partial}{\partial x_i}\right|_p=X_i|_p,\qquad i=1,\ldots,k.
\]
por \hyperref[blk:047506]{Lema} 3. Nos enfocamos en la rebanada

\[
R=\{(r_1,\ldots,r_n)\in C_{\epsilon}^{n} (0):r_1=\cdots=r_k=0\}\simeq C_{\epsilon}^{n-k} (0)\subseteq \psi(V)
\]
que tiene dimensión $n-k$. 4. Armamos el mapa

\[
\sigma:(-\varepsilon,\varepsilon)^k\times C_{\epsilon}^{n-k}\subseteq \mathbb{R}^{n}  \to M
\]
por

\[
\sigma(t_1,\ldots,t_k,r^{k+1},\ldots,r^{n}  )=\theta^1_{t_1}\circ\cdots\circ\theta^k_{t_k}(\psi^{-1}(0,\ldots 0,r^{k+1},\ldots,r^{n}  )).
\]
Notar que esto este bien definido por que al ser carta cubica $\psi ^{-1}(0,\ldots,0,r^{k+1},\ldots r^{n})$ cae en $V$ 5. Como en la vez anterior, vamos a ver que $\sigma$ es un difeomorfismo entre un abierto de $\mathbb{R}^n$ y un abierto de $M$ que contiene a $p$, mostrando que

\[
(d\sigma)_0\left(\left.\frac{\partial}{\partial r^i}\right|_0\right)=\begin{cases} X_i|_p, & i=1,\ldots,k,\\[4pt] \left.\dfrac{\partial}{\partial x^i}\right|_p, & i>k.\end{cases}
\]
\begin{enumerate}[start=6]
\item En efecto, si $i\leq k$,
\end{enumerate}
\[
\begin{aligned}(d\sigma)_0\left(\left.\frac{\partial}{\partial t^i}\right|_0\right)&=\left.\frac{d}{ds}\right|_{s=0}\sigma(se_i)\\&=\left.\frac{d}{ds}\right|_{s=0}\sigma(0,\ldots,0,s,0,\ldots,0)\\&=\left.\frac{d}{ds}\right|_{s=0}\left(\theta^{1}_{0}\circ\cdots\circ\theta^{i}_{s}\circ\cdots\circ\theta^{k}_{0}\right)\left(\psi^{-1}(0,\ldots,0)\right)\\&=\left.\frac{d}{ds}\right|_{s=0}\theta^{i}_{s}(p)\\&=X_i|_p\end{aligned}
\]
\begin{enumerate}[start=7]
\item Si $i>k$,
\end{enumerate}
\[
\begin{aligned}(d\sigma)_0\left(\left.\frac{\partial}{\partial r^i}\right|_0\right)&=\left.\frac{d}{ds}\right|_{s=0}\sigma(se_i)\\&=\left.\frac{d}{ds}\right|_{s=0}\sigma(0,\ldots,0,s,0,\ldots,0)\\&=\left.\frac{d}{ds}\right|_{s=0}(\theta^{1}_{0}\circ\cdots\circ\theta^{k}_{0})(\psi^{-1}(0,\ldots,0,s,0,\ldots,0)\\&=\left.\frac{d}{ds}\right|_{s=0}\psi^{-1}(0,\ldots,0,s,0,\ldots,0)\\&=(d\psi^{-1})_0\left(\left.\frac{\partial}{\partial r^i}\right|_0\right)\\&=\left.\frac{\partial}{\partial x^i}\right|_p\end{aligned}
\]
este ultimo igual sale rapido escribiendo a $(d\psi^{-1})_0\left(\left.\frac{\partial}{\partial r^i}\right|_0\right)$ en coordenadas del $T_{\phi ^{-1}(0)}M$ 8. Ahora notamos que por paso 2.

\[
\left\{ X_{1}|_{p},\ldots,X_{k}|_{p},\frac{\partial}{\partial x^{k+1}}\bigg|_{p},\ldots, \frac{\partial}{\partial x^{n}}|_{p}\right\}=\left\{ \frac{\partial}{\partial x^{n}}\bigg|_{p},\ldots, \frac{\partial}{\partial x^{k}}\bigg|_{p} , \frac{\partial}{\partial x^{k+1}}|_{p}\ldots\frac{\partial}{\partial x^{n}}\bigg|_{p} \right\}
\]
\begin{enumerate}[start=9]
\item Entonces $(d\sigma)_{0}$ es isomorfismo lineal por que manda base $\{ \frac{\partial}{\partial r_{i}} \}_{i<n}$ (notar aca llamo $t_{i}:=r^{i}$ es solo notacion para cordenadas en $\mathbb{R}$) en una base $\{ \frac{\partial}{\partial x^{n}}|_{p},\ldots, \frac{\partial}{\partial x^{k}}|_{p} , \frac{\partial}{\partial x^{k+1}}|_{p}\ldots\frac{\partial}{\partial x^{n}}|_{p} \}$
\item Estamos entonces en condiciones de aplicar el Teorema de la Función Inversa. Existe un abierto $\widetilde U\subseteq\mathbb{R}^n$ que contiene a $0$ y un abierto $U\subseteq M$ que contiene a $p$ tales que
\end{enumerate}
\[
\sigma|_{\widetilde U}:\widetilde U\longrightarrow U
\]
es un difeomorfismo. 11. Como $\sigma$ es un difeomorfismo entre $\widetilde U$ y $U$, llamamos $\sigma^{-1}=\varphi$, de modo que $(U,\varphi=y^{1},\ldots,y^{i})$ es una carta coordenada. 12. Recordamos nuestro objetivo es probar que

\[
X_i|_U=\frac{\partial}{\partial y^i},\qquad i=1,\ldots,k.
\]
\begin{enumerate}[start=13]
\item Equivalentemente, queremos ver que para todo $(r^1,\ldots,r^n)\in\widetilde U$, e $1\leq i\leq k$
\end{enumerate}
\[
(d\sigma)_{(r^1,\ldots,r^n)}\left(\left.\frac{\partial}{\partial r^i}\right|_{(r^1,\ldots,r^n)}\right)=(X_i)_{\sigma(r^1,\ldots,r^n)}.
\]
por que si esto es cierto usamos que

\[
(d\sigma)_{(r^1,\ldots,r^n)}\left(\left.\frac{\partial}{\partial r^i}\right|_{(r^1,\ldots,r^n)}\right)=(d\varphi ^{-1})_{(r^1,\ldots,r^n)}\left(\left.\frac{\partial}{\partial r^i}\right|_{(r^1,\ldots,r^n)}\right)=\frac{\partial}{\partial y^{i} }\bigg|_{\varphi ^{-1}(r)=\sigma(r)}
\]
\begin{enumerate}[start=14]
\item Primero notemos que, como
\end{enumerate}
\[
[X_i,X_j]=0
\]
los flujos conmutan localmente por \hyperref[blk:26db5d]{Criterio de conmutación de flujos (Hulet 9.15)} 15. Fijemos $1\leq i\leq k$. Entonces

\[
\begin{aligned}(d\sigma)_{(r^1,\ldots,r^n)}\left(\left.\frac{\partial}{\partial r^i}\right|_{(r^1,\ldots,r^n)}\right)&=\left.\frac{d}{ds}\right|_{s=0}\sigma(r^1,\ldots,r^i+s,\ldots,r^n)\\&=\left.\frac{d}{ds}\right|_{s=0}\left(\theta^{1} _{r^1}\circ\cdots\circ\theta^{i} _{r^i+s}\circ\cdots\circ\theta_{k,r^k}\right)\left(\psi^{-1}(0,\ldots,0,r^{k+1},\ldots,r^n)\right)\\&=\left.\frac{d}{ds}\right|_{s=0}\left(\theta^{1} _{r^1}\circ\cdots\circ\theta^{i} _{r^i}\circ\theta^{i} _{s}\circ\cdots\circ\theta^{k} _{r^k}\right)\left(\psi^{-1}(0,\ldots,0,r^{k+1},\ldots,r^n)\right)\\&=\left.\frac{d}{ds}\right|_{s=0}\theta^{i} _{s}\left(\sigma(r^1,\ldots,r^n)\right)\\&=(X_i)_{\sigma(r^1,\ldots,r^n)}.\end{aligned}
\]
donde usamos que los flujos conmutan

\bookqed
\end{bookproof}
